\begin{textsample}{POS Dim 3 – to – Score 17.00 – to/to\_inf\_23.txt}  \label{ex:f3_pos_095}
Elementos da \textbf{rotina} O trabalho com \textbf{crianças} \textbf{pequenas} exige a estruturação de uma \textbf{rotina} de ações, que leve em consideração as necessidades de desenvolvimento das \textbf{crianças}.
Nesse sentido, é imprescindível que o professor estruture a intencionalidade pedagógica, dentro de cada elemento, em uma \textbf{rotina} pré-estabelecida.
Elementos da \textbf{rotina} Falar sobre \textbf{rotina} na Educação \textbf{Infantil} pode ser perigoso, pois se considera que a determinação de tempos e ações engessa o processo pedagógico, e dificulta as vivências de variadas experiências por parte das \textbf{crianças}.
Entretanto, é importante destacar que a organização do tempo favorece o desenvolvimento \textbf{infantil}, uma vez que a dinâmica da Unidade torna mais segura e acessível para todas as \textbf{crianças}.
Estabelecer \textbf{rotinas} não pode ser sinônimo de monotonia, nem de momentos e ações educativas frias e descontextualizadas.
A \textbf{rotina} deve ser pensada de forma flexível, baseando -se sempre na \textbf{escuta} sensível das \textbf{crianças} por parte dos profissionais da educação.
Ela deve acontecer na perspectiva de participação das \textbf{crianças}, nos processos de construção e desconstrução das experiências e vivências no dia a dia.
A proposta de trabalho da \textbf{rotina} pode ser estabelecida em modalidades organizativas, ou seja, em atividades permanentes, sequências didáticas e projetos didáticos.
A primeira deve ocorrer todos os dias, ou em determinados momentos da semana.
A segunda é vinculada a uma proposta de trabalho sequenciado, e que deve partir do interesse das \textbf{crianças}.
E a última, relaciona -se ao desenvolvimento de um projeto que deve ocorrer de acordo com o planejamento do mesmo.
Partindo do princípio das modalidades organizativas, apresentamos a seguir a estruturação de possíveis atividades, que poderão ocorrer diariamente na Educação \textbf{Infantil}.
E sendo concebidas como atividades permanentes, devem ser pensadas com intencionalidade pedagógica, associando o cuidar e o educar, assim como os eixos estruturantes interações e \textbf{brincadeiras}. - Acolhida : O momento de chegada à Unidade Escolar pode ser de tensão e estresse para as \textbf{crianças}, pais e professores.
Para amenizar os impactos desse momento e da separação da família, é importante que a equipe pedagógica pense em diferentes estratégias de recepção e \textbf{acolhida} das \textbf{crianças}.
Com músicas, \textbf{brinquedos}, cantinhos entre outras. - Roda de conversa : Momento da \textbf{rotina} no qual as \textbf{crianças} são estimuladas a desenvolver a oralidade e expressão.
É importante que ocorra todos os dias e em diferentes momentos da \textbf{rotina}.
Deve ser planejada com variadas possibilidades e estratégias pedagógicas, tais como : caixa surpresa, notícias do dia, figuras de revistas, álbuns seriados dentre outras. - Cantos de experiência/cantinhos : É importante que seja estabelecido na \textbf{rotina} o momento de exploração dos cantos de interesse, sejam eles dentro da sala de atividade ou nos espaços externos.
Os cantos de experiências são espaços organizados do lado interno ou externo da Unidade, podem ter os mais variados nomes, devem ser planejados e preparados antes de serem disponibilizados para as \textbf{crianças}. - Leitura feita pelo professor/contação de história : Interessante que atividades de leitura de livros literários aconteçam todos os dias.
Para tanto, é importante que professores façam a leitura das obras escolhidas, antes mesmo do planejamento, incentivando a fala da \textbf{criança} acerca de suas percepções sobre a obra.
Lembrando que a leitura de história é uma forma de apresentar a obra conforme sua linguagem original, nas palavras do autor, enquanto a contação de histórias, envolve a improvisação, a interação com as \textbf{crianças} e a possibilidade de agregar outros elementos ao enredo. - Leitura feita pela \textbf{criança} : Diariamente a \textbf{criança} precisa ter contato com os livros de literatura, bem como com outros suportes de texto ou gêneros textuais.
Importante que o professor planeje, com antecedência, a disposição dos materiais e a forma de exploração dos mesmos. - Desenho : As \textbf{crianças} precisam desenhar.
O planejamento docente, e a disponibilização de materiais com antecedência, são necessários para a realização do grafismo por parte das \textbf{crianças}. - \textbf{Brincadeira} : No \textbf{pátio} externo, em parquinho de \textbf{areia}, gramado, em área cimentada ou de terra, em casinha ou perto de uma árvore.
É importante que a \textbf{criança} saia para a área externa, todos os dias, para \textbf{brincar} com seus \textbf{colegas} de sala e com \textbf{colegas} de outras salas.
É necessário que o docente planeje esse momento pedagógico, considerando sempre as possibilidades de interação entre \textbf{crianças} de diferentes faixas etárias.
É fundamental que a utilização dos espaços seja dialogada com a \textbf{criança}, com a finalidade de combinados prévios, que não privem as relações e interações, mas que fortaleçam laços de respeito, solidariedade e cooperação quanto à atuação individual na coletividade. - \textbf{Refeições} : É recomendada que a \textbf{refeição} seja planejada e acompanhada pelo professor, visto que é nesse momento que a \textbf{criança} pode ser estimulada a adquirir \textbf{hábitos} saudáveis, tanto na alimentação, quanto na higienização das mãos, antes de \textbf{manipular} os alimentos.
O professor pode apresentar o cardápio do dia, explicando de forma bem \textbf{lúdica} sobre os alimentos que serão ofertados e de onde eles vêm.
Ressalta -se a importância de ser um momento de estímulo à autonomia. - Soninho : O repouso é fundamental para a \textbf{criança}, entretanto, é importante que se compreenda que, nem toda \textbf{criança} gosta de dormir.
Nesse caso, o planejamento de ações a serem realizadas com as \textbf{crianças} que não dormem é essencial para garantir o \textbf{equilíbrio}, entre os que dormem e os que não dormem. - Banho : A atividade \textbf{lúdica} e interativa é essencial, assim como o planejamento e o preparo docente, prevendo ações de estímulo da autonomia da \textbf{criança}.
É um importante momento da \textbf{rotina}, no qual se dará atenção individualizada à \textbf{criança}, e nesse caso, é uma excelente oportunidade de aproximação e estímulo ao desenvolvimento da oralidade.
A \textbf{rotina} deve valorizar as relações cotidianas entre todos os envolvidos no processo pedagógico.
Aqui foram apresentadas algumas possibilidades de atividades de \textbf{rotina}, a serem desenvolvidas com \textbf{crianças} da Educação \textbf{Infantil} e não devem ser consideradas como as únicas vivências a serem oportunizadas às \textbf{crianças}.
Cabe à equipe pedagógica, analisar cada atividade dentro de sua realidade e cultura local.

% matched lemmas: acolher, areia, brincadeira, brincar, brinquedo, colega, criança, equilíbrio, escuta, hábito, infantil, lúdico, manipular, pequeno, pátio, refeição, rotina
\end{textsample}
